\begin{ExerciseBlock}
\ExerciseHeading{03}{Überholvorgang zweier Fahrzeuge}{linear_kinematics_overtaking_vehicles}
Ein LKW mit einer Länge $L_\text{LKW}=15$m und einer konstanten Geschwindigkeit $v_\text{LKW} = 60$km/h wird von einem PKW mit eine Länge $L_\text{PKW} = 4$m und einer Geschwindigkeit von $v_\text{PKW} = 80$km/h überholt. dabei soll der Überholvorgang beginnenm wenn die Fahrzeuge einen Abstand von $a_v = 30$m haben, und bei einem Abstand von $a_n = 40$m enden.

\begin{enumerate}
	\item Wie lange dauert der Überholvorgang?

	\item Wie lange dauert der Überholvorgang, wenn der PKW von Beginn des Überholvorgangs mit $a_\text{PKW} = 1 \, m/s^2$ beschleunigt.

\end{enumerate}

Hinweis: Fertigen Sie, bevor Sie rechnen, eine Skizze des Überholvorgangs an.
\end{ExerciseBlock}
\UseExerciseSolution{linear_kinematics_overtaking_vehicles}
